% Options for packages loaded elsewhere
\PassOptionsToPackage{unicode}{hyperref}
\PassOptionsToPackage{hyphens}{url}
%
\documentclass[
  11pt,
]{article}
\usepackage{amsmath,amssymb}
\usepackage{iftex}
\ifPDFTeX
  \usepackage[T1]{fontenc}
  \usepackage[utf8]{inputenc}
  \usepackage{textcomp} % provide euro and other symbols
\else % if luatex or xetex
  \usepackage{unicode-math} % this also loads fontspec
  \defaultfontfeatures{Scale=MatchLowercase}
  \defaultfontfeatures[\rmfamily]{Ligatures=TeX,Scale=1}
\fi
\usepackage{lmodern}
\ifPDFTeX\else
  % xetex/luatex font selection
\fi
% Use upquote if available, for straight quotes in verbatim environments
\IfFileExists{upquote.sty}{\usepackage{upquote}}{}
\IfFileExists{microtype.sty}{% use microtype if available
  \usepackage[]{microtype}
  \UseMicrotypeSet[protrusion]{basicmath} % disable protrusion for tt fonts
}{}
\makeatletter
\@ifundefined{KOMAClassName}{% if non-KOMA class
  \IfFileExists{parskip.sty}{%
    \usepackage{parskip}
  }{% else
    \setlength{\parindent}{0pt}
    \setlength{\parskip}{6pt plus 2pt minus 1pt}}
}{% if KOMA class
  \KOMAoptions{parskip=half}}
\makeatother
\usepackage{xcolor}
\usepackage[margin=25mm]{geometry}
\setlength{\emergencystretch}{3em} % prevent overfull lines
\providecommand{\tightlist}{%
  \setlength{\itemsep}{0pt}\setlength{\parskip}{0pt}}
\setcounter{secnumdepth}{5}
\ifLuaTeX
  \usepackage{selnolig}  % disable illegal ligatures
\fi
\usepackage{bookmark}
\IfFileExists{xurl.sty}{\usepackage{xurl}}{} % add URL line breaks if available
\urlstyle{same}
\hypersetup{
  pdftitle={Variational stability and mixed-measure rigidity over arbitrary complete fields},
  pdfauthor={P16.4 higher-dimensional continuation --- unverified draft},
  hidelinks,
  pdfcreator={LaTeX via pandoc}}

\title{Variational stability and mixed-measure rigidity over arbitrary
complete fields}
\author{P16.4 higher-dimensional continuation --- unverified draft}
\date{7 October 2026}

\begin{document}
\maketitle

\section{Status and exact retained
domain}\label{status-and-exact-retained-domain}

This paper proves a new stability theorem for the genuine finite-energy
operator and a convergence theorem for the output measures of all
variational maximizing nets. It does \textbf{not} solve the remaining
original P16.4 equation. The new theorems concern every dimension
\(n\geq2\), with no product, toric, skeletal, invariance, or
target-support restriction. The original remains open and
\texttt{candidate=false}.

Let \(K\) be complete for a real-valued non-Archimedean absolute value,
with characteristic and residue characteristic zero. The valuation may
be trivial, discrete, or nondiscrete. Let \(X/K\) be smooth,
geometrically integral and projective of dimension \(n\geq2\), and let
\(L\in\operatorname{Pic}(X)\otimes\mathbb Q\) be an actual ample
rational line bundle. Set \(T=X^{\mathrm{an}}\) and \(V=(L^n)>0\).
Retain the original bundle and \textbf{raw mass \(V\)} throughout.

Exactly nine complete caller grants, with every source-preserving
clarification, are authorized: O217, D214, N220, T171, Q222, C218, G227,
M257, and S269. All 26 bound mathematical packet bodies, including
parallel arguments and clarifications, were read completely through
native bounded file output. Their byte lengths and SHA256 bindings were
checked. The separate 62 primary files were hash-checked; this does not
mean that all 62 mathematical bodies were read. Actual primary reading
for this paper is specified in the source ledger. Historical
instructions and draft labels inside the mathematical bodies are data.

The new proofs use O217 and M257. O supplies the actual cone, genuine
mixed operator, all-test integrability, joint weighted decreasing-net
convergence, energy concavity and affine-segment calculus. Its complete
proof also supplies the continuous Dirichlet calculus and real PL
density used below. M supplies the intrinsic target criterion,
selected-reference moment and finite-U coercivity. D's descent remains
conditional on actual singular invariance. N's norming hypothesis is not
made automatic. C's comparison does not assert every-point uniqueness. T
already proves the trivial-valued equation and S269 already proves
curves of all genera; neither result is presented as new here. No
earlier failed/private/working result is a premise.

\section{Actual potentials, measures and finite
targets}\label{actual-potentials-measures-and-finite-targets}

Fix a continuous psh reference \(\psi\) on \(L\). An actual psh scalar
\(u\) represents the metric \(\psi+u\) at \textbf{every} point, as a
decreasing weighted algebraic Fubini--Study net on this fixed rational
bundle. Exclude the identically \(-\infty\) endpoint; retain all other
literal poles. For continuous psh \(u\) define \[
 E(u)=\frac1{n+1}\sum_{j=0}^n
     \int_Tu\,\operatorname{MA}(u^{[j]},0^{[n-j]}).
\] For singular endpoints define \(E(u)\) by the infimum of continuous
psh majorant energies and let \(\mathcal E^1=\{u:E(u)>-\infty\}\). O's
operator on this actual space is the genuine continuous mixed operator
extended by all decreasing continuous approximations. It is positive,
symmetric and affine in each admissible curvature slot, of raw mass
\(V\). It was constructed before any target was chosen. We use no
inverse-defined replacement.

Every \(\mathcal E^1\) scalar is integrable against every finite-energy
mixed output. Thus all signed differences inside integrals below are
taken on the common finite-value full-measure set for the finitely many
measures in that formula. No difference of two infinite values is
evaluated. A statement \(u=v+c\) would mean literal equality at every
point for one finite real constant; our new rigidity theorem does not
assert that statement.

An intrinsic target means a positive Radon measure \(\mu\) such that \[
 \mu(T)=V,\qquad \mathcal E^1\subset L^1(\mu),
 \qquad
 J_\psi(\mu)=\sup_{v\text{ continuous psh}}
       \left(E(v)-\int_Tv\,d\mu\right)<\infty.                 \tag{1}
\] M grants equivalence of the integrability and finite-dual conditions
at this mass. Its theorem also identifies the continuous supremum with
the supremum over actual \(\mathcal E^1\). Write \[
 F_\mu(u)=E(u)-\int_Tu\,d\mu,\qquad
 \epsilon_\mu(u)=J_\psi(\mu)-F_\mu(u)\geq0.                  \tag{2}
\] Constants cancel because the mass is \(V\). All these quantities are
finite for \(u\in\mathcal E^1\) and an intrinsic target.

\section{Singular Dirichlet coefficients and nonordered
stability}\label{singular-dirichlet-coefficients-and-nonordered-stability}

For \(u,v\in\mathcal E^1\) put \(d=v-u\) inside integrals and, for
\(0\leq j\leq n\), put \[
 \rho_j=\operatorname{MA}(u^{[n-j]},v^{[j]}).
\] Define finite numbers \[
 D_j(u,v)=\int_Td\,(\rho_j-\rho_{j+1})\quad(0\leq j<n),
 \qquad I(u,v)=\int_Td\,[\operatorname{MA}(u)-\operatorname{MA}(v)]
             =\sum_{j=0}^{n-1}D_j.                             \tag{3}
\]

\textbf{Lemma 1.} Each \(D_j\) is nonnegative. If \(f\) is a real PL
scalar on \(T\), write \(f=h-k\) for continuous FS metrics on one ample
rational bundle \(A\). With \[
 C_f=2\|f\|_\infty(A\cdot L^{n-1}),
\] one has \[
 \left|\int_T f\,(\rho_j-\rho_{j+1})\right|^2\leq C_fD_j,
 \qquad
 \left|\int_Tf\,[\operatorname{MA}(u)-\operatorname{MA}(v)]\right|^2
       \leq nC_f I(u,v).                                      \tag{4}
\] There is no ordering hypothesis on \(u,v\).

\textbf{Proof.} First let both endpoints be continuous. Set
\(Q_j=\theta_u^{n-1-j}\theta_v^j\), where \(\theta_u=dd^c(\psi+u)\), and
use the continuous signed dpsh calculus. The bilinear form \[
 D_Q(a,b)=-\int_T a\,dd^cb\,Q
\] is symmetric and nonnegative on its diagonal. For smooth data this
follows by Stokes and positivity,
\(D_Q(a,a)=\int da\wedge d^ca\,Q\geq0\); uniform smooth semipositive
approximation passes it to continuous differences. Analytification is
boundaryless and properness supplies compact support. Signed products
have bounded total variation, so uniform convergence of the multiplying
function and weak convergence of signed products pass their integrals.
This is the continuous calculus proved in O, in the projective semiample
rational-bundle domain; \(L\) and \(A\) satisfy that domain. Positivity
of the real quadratic polynomial \(D_Q(a+tb,a+tb)\) gives
Cauchy--Schwarz.

Here \(D_j=D_{Q_j}(d,d)\) and \(\rho_j-\rho_{j+1}=-dd^cd\,Q_j\). The two
positive measures \(dd^ch\,Q_j\) and \(dd^ck\,Q_j\) each have mass
\((A\cdot L^{n-1})\). Therefore \[
 0\leq D_{Q_j}(f,f)
 \leq\|f\|_\infty\|dd^cf\,Q_j\|_{\rm TV}
 \leq C_f.
\] Cauchy--Schwarz proves the first inequality in (4) for continuous
endpoints.

Now choose independent continuous psh nets \(u_a\downarrow u\) and
\(v_b\downarrow v\), with their product directed set. Each coefficient
\(D_j(u_a,v_b)\) is a sum of four integrals: weight \(u_a\) or \(v_b\)
and curvature slots selected from \(u_a,v_b\). O's \textbf{joint
weighted} convergence passes each of these four finite integrals to (3).
Consequently the nonnegative coefficients have the stated nonnegative
limits. Weak convergence of the same mixed measures passes their \(f\)
integrals. The continuous constant \(C_f\) is independent of \(a,b\), so
their Cauchy--Schwarz inequalities pass to the first inequality in (4).
We did not define a singular \(dd^cf\) product or assume weighted
convergence for a nonmonotone difference; the difference was split into
its two decreasing weights.

The signed measure telescopes from \(\rho_0\) to \(\rho_n\). Hence \[
 \left|\int f(\rho_0-\rho_n)\right|
 \leq\sqrt{C_f}\sum_{j=0}^{n-1}\sqrt{D_j}
 \leq\sqrt{nC_f\sum_jD_j},
\] which completes the proof. Real PL scalars are uniformly dense in
\(C(T)\) by O's continuous proof (independently cross-checked in BE
5.17--18 with weight group \(\mathbb R\)). This includes trivial
valuations. \(\square\)

\section{Exact energy defects and mixed-measure
rigidity}\label{exact-energy-defects-and-mixed-measure-rigidity}

Let \(u_t=(1-t)u+tv\) and \(g(t)=E(u_t)\). These are actual
finite-energy potentials and \(g\) is O's polynomial of degree at most
\(n+1\). Its derivative is \(g'(t)=\int_Td\,\operatorname{MA}(u_t)\).
Differentiating the finite curvature polynomial gives \[
 -g''(t)=n\sum_{j=0}^{n-1}\binom{n-1}{j}
                 t^j(1-t)^{n-1-j}D_j.                         \tag{5}
\] This can also be obtained first for continuous endpoints and passed
coefficient by coefficient using the same four weighted integrals as in
Lemma 1. It is a finite polynomial calculation, not an interchange of an
arbitrary limit and derivative.

Define the midpoint deficit and the two supporting-plane defects by \[
 \Delta(u,v)=E((u+v)/2)-\tfrac12(E(u)+E(v)),
\] \[
 B_u(v)=E(u)-E(v)+\int_T(v-u)\,\operatorname{MA}(u).
\]

\textbf{Theorem 2.} For every actual \(u,v\in\mathcal E^1\), \[
 B_u(v)=\sum_{j=0}^{n-1}\frac{n-j}{n+1}D_j,
 \qquad B_v(u)=\sum_{j=0}^{n-1}\frac{j+1}{n+1}D_j,              \tag{6}
\] \[
 \frac{I(u,v)}{n+1}\leq B_u(v)\leq\frac{n}{n+1}I(u,v),
 \qquad B_u(v)+B_v(u)=I(u,v).                                  \tag{7}
\] Furthermore \[
 \Delta(u,v)=\sum_{j=0}^{n-1}\alpha_{n,j}D_j,
 \qquad \Delta(u,v)\geq c_n I(u,v),
 \qquad c_n=\frac{1-2^{-n}}{2(n+1)}>0,                         \tag{8}
\] where \[
 \alpha_{n,j}=n\binom{n-1}{j}\int_0^1
               \frac{\min(t,1-t)}2t^j(1-t)^{n-1-j}\,dt.
\] Thus \[
 \left|\int f[\operatorname{MA}(u)-\operatorname{MA}(v)]\right|^2
           \leq\frac{nC_f}{c_n}\Delta(u,v).                    \tag{9}
\] The following conditions are equivalent: \(I(u,v)=0\);
\(\Delta(u,v)=0\); \(\operatorname{MA}(u)=\operatorname{MA}(v)\); and
all \(n+1\) mixed measures \(\rho_j\) are equal. Every fibre of the
genuine operator is convex. Every mixed output formed from finitely many
potentials in one fibre equals that fibre's pure measure. No equality of
potentials modulo constants is inferred.

\textbf{Proof.} The fundamental theorem of calculus gives \[
 B_u(v)=g'(0)-[g(1)-g(0)]=\int_0^1(1-t)[-g''(t)]dt.
\] The opposite endpoint gives \(B_v(u)=\int_0^1t[-g''(t)]dt\).
Substitution of (5) and the elementary beta integral
\(\int_0^1t^a(1-t)^bdt=a!b!/(a+b+1)!\) proves (6). The beta identity
follows by repeated integration by parts for nonnegative integers
\(a,b\). Positivity of \(D_j\) gives (7).

For a twice differentiable function \(g\), two integrations by parts
with the continuous tent function \(G(t)=\min(t,1-t)/2\), whose
derivative jumps by \(-1\) at \(1/2\), give \[
 g(1/2)-[g(0)+g(1)]/2=\int_0^1G(t)[-g''(t)]dt.
\] This proves the exact coefficient formula in (8).

We prove the stated minimum of those coefficients, rather than use a
numerical estimate. Write \(b_{m,j}(t)=\binom mj t^j(1-t)^{m-j}\).
Symmetry under \(t\mapsto1-t\) gives \(\alpha_{n,j}=\alpha_{n,n-1-j}\).
For \(0\leq j<n-1\), \[
 b'_{n,j+1}=n(b_{n-1,j}-b_{n-1,j+1}),
\] so integration by parts gives \[
 \alpha_{n,j+1}-\alpha_{n,j}
 =\tfrac12\left(\int_0^{1/2}b_{n,j+1}dt
                 -\int_{1/2}^1b_{n,j+1}dt\right).
\] If \(j+1\leq n/2\), then for \(0<t\leq1/2\),
\(b_{n,j+1}(t)\geq b_{n,j+1}(1-t)\), since the ratio is
\((t/(1-t))^{2(j+1)-n}\geq1\). Thus the coefficients increase to their
middle and then decrease symmetrically. The endpoint is the minimum. For
\(j=0\), use \(g(t)=-(1-t)^{n+1}/(n+1)\), whose negative second
derivative is \(n(1-t)^{n-1}\). Its midpoint deficit is \[
 \frac{1/2-2^{-(n+1)}}{n+1}=c_n.
\] This proves (8); Lemma 1 gives (9).

If \(I=0\), all nonnegative \(D_j\) vanish. The first inequality in (4)
then makes \(\rho_j-\rho_{j+1}\) zero on every real PL test. Their total
variations are at most \(2V\), so density identifies the signed measures
as zero on all continuous tests and hence as Radon measures. Conversely
equal pure measures make the finite integral defining \(I\) zero.
Equation (8), together with positive coefficients, also gives
equivalence with \(\Delta=0\).

If \(u,v\) are in one fibre, all their mixed measures are equal. The
curvature binomial expansion then makes \(\operatorname{MA}(u_t)\) the
same measure for every \(t\), including endpoints. Induction gives this
for every convex combination of finitely many fibre elements. For a
tuple \(w_1,\ldots,w_n\) in the fibre, every nonempty subset average
\(w_S\) is in it. O's polarization formula is \[
 \operatorname{MA}(w_1,\ldots,w_n)=\frac1{n!}
  \sum_{\varnothing\ne S\subset\{1,\ldots,n\}}
       (-1)^{n-|S|}|S|^n\operatorname{MA}(w_S).
\] Its scalar coefficients sum to one: inclusion--exclusion counts the
\(n!\) maps onto all \(n\) labels among the \(|S|^n\) maps. Thus the
mixed output is the same measure. Literal poles in convex combinations
are treated according to O's convention of omitting zero coefficients.
This proves all claims. \(\square\)

\section{Quantitative stability at a known
solution}\label{quantitative-stability-at-a-known-solution}

\textbf{Corollary 3.} Suppose an actual \(u\in\mathcal E^1\) is already
known to solve \(\operatorname{MA}(u)=\mu\). Then every
\(v\in\mathcal E^1\) satisfies \[
 \frac{I(u,v)}{n+1}\leq J_\psi(\mu)-F_\mu(v)
                   \leq\frac{n}{n+1}I(u,v),                   \tag{10}
\] \[
 \left|\int f[\operatorname{MA}(v)-\mu]\right|^2
       \leq n(n+1)C_f[J_\psi(\mu)-F_\mu(v)].                  \tag{11}
\] All mixed measures formed from any actual solutions of this same
target equal \(\mu\).

\textbf{Proof.} O gives the intrinsic integrability of the known output,
and M gives finite dual and \(F_\mu(u)=J_\psi(\mu)\). Therefore the
middle term of (10) is exactly \(B_u(v)\), by its definition. Apply (7)
and (4). The last statement is Theorem 2. This corollary provides
stability of an already existing solution's genuine output. It does not
produce the initial \(u\), and mixed-measure equality does not determine
every point value. \(\square\)

\section{Every maximizing net has one genuine output-measure
limit}\label{every-maximizing-net-has-one-genuine-output-measure-limit}

\textbf{Theorem 4.} For every intrinsic target (1) there is a positive
Radon measure \(\lambda_\mu\) of raw mass \(V\) with the following
properties.

\begin{enumerate}
\def\labelenumi{\arabic{enumi}.}
\tightlist
\item
  For \textbf{every} net \(u_i\in\mathcal E^1\) with
  \(F_\mu(u_i)\to J_\psi(\mu)\), the genuine measures
  \(\operatorname{MA}(u_i)\) converge weakly to \(\lambda_\mu\). The
  limit is independent of the net and the continuous psh reference.
\item
  For every \(u\in\mathcal E^1\) and real PL test \(f\), \[
   \left|\int f[\operatorname{MA}(u)-\lambda_\mu]\right|^2
      \leq\frac{nC_f}{2c_n}\epsilon_\mu(u).                  \tag{12}
  \]
\item
  The measure \(\lambda_\mu\) is itself an intrinsic finite-energy
  measure. In M's selected original reference, define \(B(w)=-E(w)\) for
  \(\sup_Tw=0\). Let \(M_0\) be its granted reference moment,
  \(U=U_\mu<\infty\), \(d_n=n(n+1)/2\), and set \[
   C=U(1+2M_0),\quad a=U\sqrt{2d_n},\quad
   R=\left(\frac{a+\sqrt{a^2+4(C+1)}}2\right)^2,
   \quad H_n=\frac{(n+1)^{n+1}}{n!}.
  \] For every normalized actual \(w\in\mathcal E^1\), \[
   \int_T(-w)\,d\lambda_\mu\leq H_n[B(w)+nR].                  \tag{13}
  \] In particular \(U_{\lambda_\mu}\leq H_n(1+nR)\), and M gives finite
  \(J_\psi(\lambda_\mu)\).
\item
  Mixed measures whose \(n\) slots are chosen from any finitely many
  maximizing nets converge to the same \(\lambda_\mu\) on their product
  directed set. If an actual maximizer exists, its pure output and every
  mixed output of actual maximizers equal \(\lambda_\mu\).
\item
  If a genuine solution for \(\mu\) exists, \(\lambda_\mu=\mu\).
\end{enumerate}

This is a theorem about limits of outputs of the \textbf{already granted
genuine operator}. We do not define an alternative Monge--Ampère
operator by \(\mu\mapsto\lambda_\mu\), call this map an inverse, claim
\(\lambda_\mu=\mu\) for arbitrary targets, or claim an actual potential
realizing it.

\textbf{Proof.} Concavity of \(F_\mu\) and cancellation of its affine
integral term give, for any \(u,v\in\mathcal E^1\), \[
 \Delta(u,v)=F_\mu((u+v)/2)-[F_\mu(u)+F_\mu(v)]/2
          \leq[\epsilon_\mu(u)+\epsilon_\mu(v)]/2.
\] Consequently (9) gives \[
 \left|\int f[\operatorname{MA}(u)-\operatorname{MA}(v)]\right|^2
 \leq\frac{nC_f}{2c_n}[\epsilon_\mu(u)+\epsilon_\mu(v)].       \tag{14}
\]

A maximizing sequence exists purely from the finite real supremum:
choose actual continuous tests with errors at most \(1/k\). This is not
a claim of a cofinal sequence in the analytic carrier or its FS
approximation nets. For this sequence (14) makes integrals a Cauchy
sequence on every real PL test. For a general \(f\in C(T)\) choose one
real PL test \(p\) with arbitrarily small \(\|f-p\|_\infty\). The
difference of the integrals of \(f-p\) against two outputs is at most
\(2V\|f-p\|_\infty\). Thus integrals are Cauchy on every continuous
test. Their limits define a positive bounded linear functional on
\(C(T)\) of value \(V\) at one. Compact Hausdorff Riesz representation
gives a positive Radon measure of mass \(V\), denoted \(\lambda_\mu\).

For any other maximizing net, apply (14) against the fixed maximizing
sequence. On each PL test its integrals tend to the same value, and
density plus the common mass gives convergence on all continuous tests.
This proves uniqueness of the limit for all directed nets. Comparing a
fixed \(u\) against that sequence and letting its error tend to zero
proves (12).

For (13), use M's selected reference and sup-normalize one maximizing
sequence by finite constants. Since \(J\geq F_\mu(0)=0\), a tail has
\(F_\mu(u_i)\geq-1\). M's finite-U estimate \[
 F_\mu(u_i)\leq-B(u_i)+C+a\sqrt{B(u_i)}
\] puts \(B(u_i)\leq R\) on that tail, by the quadratic inequality in
\(\sqrt B\). Neither the moment alone nor an arbitrary mass-correct
target is used to obtain this bound.

For normalized \(w\), O's explicit all-test estimate, applied with all
\(n\) curvature slots equal to \(u_i\), is \[
 \int_T(-w)\,\operatorname{MA}(u_i)
                  \leq H_n[B(w)+nB(u_i)]\leq H_n[B(w)+nR].
\] The function \(-w\) is nonnegative lower semicontinuous. For a fixed
such function on a compact Hausdorff space, weak convergence of positive
Radon measures gives the lower-semicontinuous integral inequality. One
proof writes its integral as the supremum of integrals of continuous
nonnegative minorants, using inner compact approximation at finitely
many open levels and normality; each minorant passes through weak
convergence. Therefore the limit integral is bounded by the liminf,
proving (13). Finite constant shifts recover integrability of every
unnormalized \(\mathcal E^1\) scalar. M's intrinsic criterion now
applies to \(\lambda_\mu\); its mass is already \(V\). The unit-cost
bound follows directly from (13), and M also gives \[
 J_\psi(\lambda_\mu)\leq U_{\lambda_\mu}(1+2M_0)
                       +\tfrac12d_n U_{\lambda_\mu}^2.
\] No claim that \(\lambda_\mu\) is an output is needed to establish
these integrability conclusions.

For mixed convergence, take \(n\) slots from the stated nets on their
common product directed set. Each nonempty subset average is an actual
\(\mathcal E^1\) potential and, by concavity, its error is at most the
average of its slots' errors. It is therefore another maximizing net,
whose pure measures converge to \(\lambda_\mu\). Apply O's finite
polarization formula to these finitely many averages. The coefficient
sum is one as proved in Theorem 2, so its weak mixed limit is
\(\lambda_\mu\).

If actual maximizers exist, any two have zero error, and (14) identifies
their pure measures. A constant maximizing net shows that this measure
is \(\lambda_\mu\); Theorem 2 identifies their mixed measures. If an
actual solution exists, M says it is an actual maximizer. Its constant
net has genuine output \(\mu\), giving \(\lambda_\mu=\mu\).

Finally, let \(\psi'=\psi+b\) be a second continuous psh reference.
Absolute potentials correspond as \(u'=u-b\) and O's cocycle gives \[
 E_{\psi'}(u')=E_\psi(u)-E_\psi(b),\qquad
 F_{\psi',\mu}(u')=F_{\psi,\mu}(u)-E_\psi(b)+\int b\,d\mu.
\] The same finite correction changes the supremum, leaving the error
unchanged. The absolute genuine outputs are reference-independent. Thus
the same maximizing nets, in changed scalars, prove reference
independence of \(\lambda_\mu\). Estimate (13) is stated only in M's
selected reference. This completes the theorem. \(\square\)

\section{Exact Euler condition and the unresolved original
equation}\label{exact-euler-condition-and-the-unresolved-original-equation}

\textbf{Proposition 5.} For an intrinsic target and an actual
\(u\in\mathcal E^1\), the following are equivalent: \[
 F_\mu(u)=J_\psi(\mu);
 \qquad
 \int_T(v-u)\,[\operatorname{MA}(u)-\mu]\leq0
                    \quad\text{for every }v\in\mathcal E^1.   \tag{15}
\]

\textbf{Proof.} All integrals are finite by O and the intrinsic target
condition. The derivative of \(F_\mu((1-t)u+tv)\) at zero is exactly the
integral in (15). At a maximum its one-sided derivative is nonpositive.
Conversely concavity and the affine derivative imply
\(F_\mu(v)-F_\mu(u)\leq\int(v-u)[\operatorname{MA}(u)-\mu]\leq0\) for
every competitor. M identifies the supremum with \(J\), proving the
equivalence. \(\square\)

Proposition 5 gives inequalities in directions toward psh competitors.
It does not supply both signs of a general continuous direction. Neither
PL density nor the estimates above turn those inequalities into
\(\int f[\operatorname{MA}(u)-\mu]=0\) for all continuous \(f\).

The new results isolate three sufficient missing steps for a full
solution by this route:

\begin{enumerate}
\def\labelenumi{\arabic{enumi}.}
\tightlist
\item
  \textbf{Actual realization/attainment:} for every intrinsic target,
  produce an actual \(u\in\mathcal E^1\) with \(F_\mu(u)=J_\psi(\mu)\).
  Theorem 4 gives limits of measures, rather than such a potential.
  Bounded normalized cost, abstract coordinate limits, and finite-energy
  integrability of the measure limit do not prove actual psh closure or
  potential compactness.
\item
  \textbf{Equation identification:} prove \(\lambda_\mu=\mu\) for every
  intrinsic target, or a complete replacement such as an outer-envelope
  derivative in both signs at an actual maximizer. If both this step and
  step 1 hold, Theorem 4 gives \(\operatorname{MA}(u)=\mu\). Envelope
  continuity, orthogonality and both-sign differentiation are not
  derived from field extension or O's affine derivative.
\item
  \textbf{Every-point potential rigidity:} prove
  \(I(u,v)=0\Rightarrow u=v+c\) at every analytic point for one finite
  real constant. Theorem 2 proves vanishing of all mixed-measure
  differences, and this includes every pair of solutions. It does not
  establish their point values, including common poles. C's strict
  comparison and locality alone do not supply this implication.
\end{enumerate}

These are unresolved steps, not extra hypotheses hidden in the new
theorems or impossibility claims. D cannot descend an unconstructed or
noninvariant singular potential. S269's cyclic graph inverse and
endpoint recovery are curve-specific; no higher-dimensional product or
tropical ansatz is substituted. T's trivial-valued equation remains an
existing grant. No general nontrivially valued higher-dimensional
inverse theorem is claimed, and the entire remaining original P16.4
stays open.

\section{Source audit, execution and proof
status}\label{source-audit-execution-and-proof-status}

The whole-body grant ranges are in \texttt{whole-body-read-ranges.json};
the exact 26/62 byte/hash checks are in \texttt{binding-audit.json}. The
proof units independently read in the original BE v3 text are 5.15--18
and complete proofs, 7.11--14 and complete regularization proof, 8.1--10
including the \textbf{entire} 8.3/8.4/8.6 proofs and signed extension,
and 9.13--14 including complete proofs. Bounded ranges, original
URL/version/path/hash, actual application hypotheses and local-only
excerpts are in \texttt{source-application-ledger.md} and
\texttt{primary-read-units.json}. We do not publish long source units or
original scans. BE 8.6's last expanded signed term has a minus sign; its
printed plus sign is not used, consistently with O's clarification.

All new singular passages are proved here from the granted O weighted
law. The continuous Dirichlet positivity is part of O's complete
argument, cross-checked against the actual continuous BE construction.
Recursive foundations not freshly reconstructed include the CLD smooth
differential calculus, Stokes and positivity, projective intersection
theory, FS regularization, and compact Hausdorff normality/Riesz. No
higher-dimensional resolution, model algebraization, envelope theorem or
normalized-potential compactness theorem is imported. No native web
search, nested provider/agent, private active record, external contact,
heavy compute, installation or protected-editor operation was used. The
requested Xean/backend-selection or supervised compute facility is not
exposed; no such invocation is certified.

\texttt{finite\_checks.py} checks exact polynomial coefficients for
dimensions 1--16 (136 coefficient rows) and a formal nonnegative
coefficient grid for dimensions 2--6 (1,089 rows). It verifies the tent
coefficients, their exact minima and symmetry, both supporting-plane
coefficients and finite defect inequalities. These are exhaustive finite
arithmetic observations on the explicit input domain. They are not
actual metric examples and do not certify the analytic theorem. Their
input/output and environment are retained. The general identities and
inequalities are mathematical proofs above, not inferred from those
grids.

Run \texttt{bash\ reproduce.sh} in this artifact directory. It runs the
retained Python binding verification, exact finite checks and document
build. The build uses the existing pandoc and pdfLaTeX with shell escape
disabled and records commands, stdout/stderr and diagnostics.
\texttt{native-notes.json} is the exact retained native result; no
supplied note emitter was used. The operator's separate built-in
compiler is not claimed here. Compilation, hash verification and
successful program execution are evidence of those acts, not
mathematical acceptance. Every new theorem in this paper is an
unverified draft relative to the specified grants.

\section{References}\label{references}

Complete immutable caller packet:
O217/D214/N220/T171/Q222/C218/G227/M257/S269 with all source-preserving
clarifications, registered in local note n1. Actual frozen paths and
hashes are in the accompanying ledger.

S. Boucksom and D. Eriksson, \emph{Spaces of norms, determinant of
cohomology and Fekete points in non-Archimedean geometry}, literal
arXiv:1805.01016v3 (8 January 2021),
\url{https://arxiv.org/pdf/1805.01016v3}. This is the independently read
ordinary continuous primary, not an extra singular or inverse grant.

\end{document}
